\documentclass[10pt,a4paper]{amsart}
\usepackage[margin=19mm]{geometry}
\usepackage{amsmath,amssymb,mathtools}
\usepackage[scr=rsfs,cal=euler]{mathalfa}
\usepackage{xfrac}
\usepackage[unicode,hidelinks,bookmarks=false]{hyperref}
\newcommand{\ie}{i.e.\ }
\newcommand{\cf}{cf.\ }
\newcommand{\resp}{respectively\ }
\newcommand{\Subsection}{\subsection}
\providecommand{\nobreakdash}{\nobreak\hbox{-}}
\setlength{\emergencystretch}{2em}
\multlinegap0pt
\usepackage[shortlabels]{enumitem}
\usepackage{mathtools}
\usepackage{booktabs}
\usepackage{tikz}
\usepackage{multirow}
\usepackage{float}
\usepackage{algorithm}\usepackage{algpseudocode}
\usepackage{amssymb}
\usepackage{natbib}
\newtheorem{theorem}{Theorem}
\newcommand{\be}{\begin{eqnarray}}
\newcommand{\ee}{\end{eqnarray}}
\title{Multiple Hypothesis Testing To Estimate The Number Of Communities in Stochastic Block Models}
\author{Chetkar Jha, Mingyao Li, Ian Barnett}
\date{Source: arXiv:2507.15471v2 (CC BY 4.0)}
\begin{document}
\maketitle

\noindent\emph{Source and licence.} Multiple Hypothesis Testing To Estimate The Number Of Communities in Stochastic Block Models. Version arXiv:2507.15471v2 (CC BY 4.0), distributed by arXiv under the Creative Commons Attribution 4.0 International licence, http://creativecommons.org/licenses/by/4.0/, as recorded in the arXiv licence metadata for this exact version and shown on the abstract page https://arxiv.org/abs/2507.15471v2. Full licence notice: \texttt{licenses/arxiv-2507.15471-CC-BY-4.0.txt}.\par\smallskip

\noindent\textit{Editorial scope.} Counted unit: the article's theorem est.TW.thm (source0.tex lines 349--357). The article's complete original proof of it is delivered with it (source0.tex lines 359--361, one \texttt{proof} environment of 3 lines). No mathematics is written, repaired or re-derived: the statement and the proof are the article's own lines, in the article's own notation, and no retained line is substituted or re-flowed.  Retained uncounted as the source-local context the proof needs: lines 345--347.  Nothing was omitted from any retained span, so the package carries no omission record.  No retained line contains a citation, so the wrapper carries no bibliography and no bibliography entry.  No retained line contains a figure environment, an image-inclusion command or drawing code, so no image is delivered and the package declares no asset dependency.  Editorial formatting only, in addition: an AMS \texttt{amsart} preamble that restates the article's own declarations and macros used by the retained lines, namely the environments \texttt{theorem} on separate counters, together with the standard packages those lines need.  The retained environments print in this document as Theorem 1 in order of appearance, whereas the article prints the same items with the numbering of its own full document.  The complete native source of the article as distributed in the arXiv e-print for this exact version is bundled unchanged as \texttt{sources/181938/source0.tex}.
\begin{enumerate}[ label =A\arabic*]
\item\label{A1} (Balancedness) : Assume that all the communities of SBM are of similar sizes (under the null). 
\end{enumerate}

\begin{theorem} (Asymptotic Null Distribution)\label{est.TW.thm}
Let $A$ be an adjacency matrix generated from a SBM($g, G_{\rho_n}$) with $K_{\star}^{(n)}$ satisfying assumption (\ref{A1}). Assume that $\hat{g}$ is a consistent estimate of $g$. Let $\{ \hat{M}^{(i)}\}_{i=1}^{K_{\star}^{(n)}}$ be scaled adjacency matrices corresponding to the adjacency matrices $\{ A_{\star}^{(i)} \}_{i=1}^{K_{\star}^{(n)}}$. Then, we show that the second largest eigenvalue of the Erd\"{o}s R\'{e}nyi graph converges to the Tracy-Widom distribution under the following conditions

\be\label{est.tW}
(n/K_{\star}^{(n)})^{2/3}(\lambda_2(\hat{M}^{(i)}) - 2 - \frac{1}{\hat{\mu}_i}) \xrightarrow[]{D} TW_1(\cdot),
\ee
\text{ when } $\hat{G}_{\rho_n}(i,i) \ge (n/K_{\star}^{(n)}))^{-2/3},  \forall i =1, \cdots, K_{\star}^{(n)}$
for $K_{\star}^{(n)} = O(log(n))$.
\end{theorem}

\begin{proof}
The proof is given in the Supplement.
\end{proof}

\end{document}
